% Artículo VII. Recuperación por dependencia, 10-09-2026.
% RESULTADO_RECUPERADO: integral activo 2249, fuente/colaboracion/
% partes_iv_v/tex/residencias/86_cosmologia_cerrada.tex (S04),
% líneas 2--44 (cartas seleccionadas), 174--248 (tensor y aceleración).
% Fuente S04 leída íntegramente, SHA-256:
% feb2582c3316c8b30829cfb0476bfa2f4e95889dabfba16fe43f2d835a226673.
% Se recuperan definición, unicidad, conservación, descomposición,
% intercambio y aceleración, con sus condiciones explícitas.
% FORMALIZACION_REUNIDA: relación con la fuente efectiva de 31;
% se distingue conservación total de conservación material separada.
% CERTIFICADO_NUEVO: expansión de las pruebas de traza, balance y signos;
% cálculo local de Ricci para las tres cartas ya seleccionadas por S04.
% No se modifica el rebote de 50 ni se determina aquí una amplitud s_0.
%
% DEPENDENCIAS DE INTEGRACIÓN:
% - vii:sec:realizacion-cartan: la métrica procede de la cotetrada realizada.
% - vii:sec:acoplamientos: G_int y el reloj t_0 del mismo estado HMT;
%   c_int=1 es una elección de unidades posterior a su construcción.
% - vii:eq:ecuacion-metrica-efectiva y vii:eq:balance-metrico-total, de 31:
%   se usan únicamente para identificar el residual sobre una solución.
% - T_b es el tensor publicado por el transductor material declarado;
%   su conservación separada es una condición expresamente enunciada.
% - El selector TRIT de la carta y t_0(x)>0 son antecedentes de la
%   especialización de Sitter, no datos cosmológicos observacionales.
% El presente fragmento no certifica la clausura del conjunto de VII.

\section{Tensor residual, descomposición de traza y balance covariante}
\label{vii:sec:tensor-residual-balance}

La realización geométrica del estado enriquecido proporciona una métrica,
mientras el transductor material expresa la fuente a la que se refiere la
lectura. El tensor residual conserva la diferencia entre ambas representaciones
en un mismo espacio tensorial. Su descomposición permite distinguir el
sector proporcional a la métrica y el sector de traza nula, y determina
la corriente de intercambio que corresponde a ese reparto.

\subsection{Datos de la coordenatización y construcción del tensor residual}
\label{vii:subsec:datos-tensor-residual}

Fijemos un estado HMT \(x\) y una coordenatización conexa \(\mathcal U\) de su
realización geométrica. Sobre \(\mathcal U\), la cotetrada no degenerada
de la sección~\ref{vii:sec:realizacion-cartan} determina una métrica lisa
\(g=g_x\) de firma \((-+++)\). Denotamos por
\(\mathring\nabla\) su conexión de Levi--Civita y escribimos
\[
 \mathsf G_{\mu\nu}[g]
 :=\operatorname{Ric}_{\mu\nu}[g]-\tfrac12R[g]g_{\mu\nu}.
\]
Así se distingue el tensor de Einstein \(\mathsf G[g]\) del acoplamiento
gravitatorio \(G_{\rm int}\). La derivada utilizada en los balances
de esta sección es \(\mathring\nabla\), coherentemente con la ecuación
métrica efectiva de la sección~\ref{vii:subsec:fuente-metrica-efectiva}.

Trabajamos en la coordenatización dimensional \(c_{\rm int}=1\). Suponemos
\(G_{\rm int}(x)>0\) y \(\Lambda_{\rm ref}(x)\) constantes sobre
\(\mathcal U\), y ponemos \(\kappa_x=8\pi G_{\rm int}(x)\).
El símbolo \(x\) identifica el estado que suministra esos lectores;
la diferenciación covariante actúa sobre los puntos de \(\mathcal U\).
La genealogía de los acoplamientos permanece en la
sección~\ref{v:ant:sec:gravity}.

Sea \(T_b\) el tensor simétrico liso expresado por el transductor material
de la coordenatización; en la lectura bariónica, es su fuente bariónica. Definimos
\begin{equation}
 \mathcal E^{\rm ref}_{\mu\nu}
 :=\kappa_x^{-1}\bigl(\mathsf G_{\mu\nu}[g]
                            +\Lambda_{\rm ref}g_{\mu\nu}\bigr),
 \qquad
 T^{\rm res}_{\mu\nu}:=\mathcal E^{\rm ref}_{\mu\nu}-T_{b,\mu\nu}.
 \label{vii:eq:residual-definicion}
\end{equation}
El tensor \(T^{\rm res}\) es el tensor residual efectivo, también
denotado \(T_{\rm osc}\) al interpretar el sector geométrico que queda
por expresar respecto de \(T_b\). La construcción conserva como
argumentos efectivos \(g_x,G_{\rm int},\Lambda_{\rm ref}\) y \(T_b\).

\begin{theorem}[Unicidad del residual y balance con la fuente material]
\label{vii:thm:residual-unicidad-balance}
Para los datos anteriores existe un único tensor simétrico
\(T^{\rm res}\) que satisface
\(\mathcal E^{\rm ref}=T_b+T^{\rm res}\). Su divergencia es
\begin{equation}
 \mathring\nabla^\mu T^{\rm res}_{\mu\nu}
 =-\mathring\nabla^\mu T_{b,\mu\nu}.
 \label{vii:eq:balance-residual-material}
\end{equation}
En particular, si la fuente material se conserva por separado,
\(\mathring\nabla^\mu T_{b,\mu\nu}=0\), también se conserva
\(T^{\rm res}\).
\end{theorem}
\begin{proof}
La diferencia de dos tensores que satisfagan la ecuación anunciada es
cero; la expresión~\eqref{vii:eq:residual-definicion} proporciona su
existencia. La identidad de Bianchi contraída y la compatibilidad métrica
de Levi--Civita dan
\[
 \mathring\nabla^\mu\mathsf G_{\mu\nu}=0,
 \qquad \mathring\nabla g=0.
\]
Como \(\kappa_x\) y \(\Lambda_{\rm ref}\) son constantes sobre la
coordenatización, \(\mathring\nabla^\mu\mathcal E^{\rm ref}_{\mu\nu}=0\).
Derivar la diferencia que define \(T^{\rm res}\) prueba
\eqref{vii:eq:balance-residual-material}, incluida la afirmación bajo
conservación material separada.
\end{proof}

La constancia espacial y temporal de los acoplamientos en la coordenatización fija
el dominio de este balance. Una colección de estados puede expresar
acoplamientos diferentes en coordenatizaciones diferentes; cada identidad anterior
se aplica a la coordenatización y al estado que la especifican.

\subsection{Descomposición canónica y corriente de intercambio}
\label{vii:subsec:traza-corriente-residual}

Definimos la función de traza y los dos tensores
\begin{equation}
 \theta:=g^{\mu\nu}T^{\rm res}_{\mu\nu},\qquad
 T^{\Lambda}_{\mu\nu}:=\tfrac14\theta g_{\mu\nu},\qquad
 T^{0}_{\mu\nu}:=T^{\rm res}_{\mu\nu}-T^{\Lambda}_{\mu\nu}.
 \label{vii:eq:descomposicion-traza-residual}
\end{equation}
El superíndice \(\Lambda\) identifica la componente proporcional a la
métrica; \(\Lambda_{\rm ref}\) es el parámetro de referencia de
\eqref{vii:eq:residual-definicion}. Sus funciones son distintas.

\begin{proposition}[Proyectores de traza y unicidad de la descomposición]
\label{vii:prop:proyectores-traza-residual}
Sobre los tensores simétricos covariantes de orden dos, las aplicaciones
\[
 P_{\rm tr}(S)=\tfrac14\operatorname{tr}_g(S)g,
 \qquad P_0(S)=S-P_{\rm tr}(S)
\]
satisfacen
\begin{equation}
 P_{\rm tr}^2=P_{\rm tr},\quad P_0^2=P_0,\quad
 P_{\rm tr}P_0=P_0P_{\rm tr}=0,\quad P_{\rm tr}+P_0=I.
 \label{vii:eq:proyectores-traza-residual}
\end{equation}
En particular, \(T^{\rm res}=T^\Lambda+T^0\),
\(\operatorname{tr}_gT^0=0\), y esta descomposición es única.
\end{proposition}
\begin{proof}
En dimensión cuatro, \(g^{\mu\nu}g_{\mu\nu}=4\). Por tanto,
\(\operatorname{tr}_g(P_{\rm tr}S)=\operatorname{tr}_gS\), lo que
prueba la idempotencia de \(P_{\rm tr}\), la anulación de la traza
de \(P_0S\) y, por expansión, las otras identidades.
Si \(S=fg+Q\) con \(\operatorname{tr}_gQ=0\), tomar la traza
da \(f=\operatorname{tr}_gS/4\). Queda determinado también
\(Q=S-fg\), y con ello la unicidad.
\end{proof}

\begin{theorem}[Intercambio covariante entre los dos sectores]
\label{vii:thm:intercambio-traza-residual}
En el dominio de \eqref{vii:eq:residual-definicion}, sea
\(b_\nu:=\mathring\nabla^\mu T_{b,\mu\nu}\). La corriente
de intercambio de traza es la uno-forma
\begin{equation}
 \mathcal J^{\rm tr}_\nu
 :=\tfrac14\partial_\nu\theta
 =\mathring\nabla^\mu T^\Lambda_{\mu\nu}.
 \label{vii:eq:corriente-intercambio-traza}
\end{equation}
Los balances completos son
\begin{equation}
 \mathring\nabla^\mu T^0_{\mu\nu}
       =-b_\nu-\mathcal J^{\rm tr}_\nu,
 \qquad
 \mathring\nabla^\mu(T^\Lambda+T^0)_{\mu\nu}=-b_\nu.
 \label{vii:eq:balances-sectores-residuales}
\end{equation}
Si \(b=0\), ambos sectores intercambian corrientes exactamente opuestas.
En ese caso se conservan por separado si y sólo si \(\theta\) es
constante en la coordenatización conexa.
\end{theorem}
\begin{proof}
La compatibilidad métrica permite calcular
\[
 \mathring\nabla^\mu\bigl(\tfrac14\theta g_{\mu\nu}\bigr)
 =\tfrac14g^{\mu\alpha}(\partial_\alpha\theta)g_{\mu\nu}
 =\tfrac14\partial_\nu\theta.
\]
Restar esta igualdad a \eqref{vii:eq:balance-residual-material} produce
\eqref{vii:eq:balances-sectores-residuales}. Con \(b=0\), la
conservación separada equivale a \(\mathrm d\theta=0\); una función
lisa con diferencial nulo es constante en cada componente conexa.
\end{proof}

La uno-forma \(\mathcal J^{\rm tr}\) mide el intercambio del reparto
tensorial. Conserva un tipo distinto del covector de transporte de la
sección~\ref{vii:sec:variacion-memoria} y de la corriente de espín
\(\sigma_{IJ}\in\Omega^3(\mathcal U;\Lambda^2E^*)\) de
\eqref{vii:eq:corriente-variacional}.

\subsection{Enlace con la fuente efectiva de Cartan--Holst}
\label{vii:subsec:residual-fuente-cartan}

Sobre una solución de \eqref{vii:eq:ecuacion-metrica-efectiva}, en las
mismas unidades \(c=1\), el residual tiene la expresión
\begin{equation}
 T^{\rm res}
 =T^{\rm phys,LC}+U^{\rm phys,spin}-T_b
   +\frac{\Lambda_{\rm ref}-\Lambda_{\rm int}}{\kappa_x}\,g.
 \label{vii:eq:residual-composicion-cartan}
\end{equation}
En efecto, sustituir
\(\mathsf G+\Lambda_{\rm int}g
  =\kappa_x(T^{\rm phys,LC}+U^{\rm phys,spin})\)
en \eqref{vii:eq:residual-definicion} da la igualdad término a término.
Cuando el transductor material expresa precisamente
\(T_b=T^{\rm phys,LC}\) y \(\Lambda_{\rm ref}=\Lambda_{\rm int}\),
se obtiene \(T^{\rm res}=U^{\rm phys,spin}\).
La condición de conservación separada de \(T_b\) del
teorema~\ref{vii:thm:residual-unicidad-balance} conserva su función
propia: \eqref{vii:eq:balance-metrico-total} establece el balance de la
fuente total, y el reparto entre sus componentes depende de la acción
material utilizada.

\subsection{Sector de traza y aceleración de las coordenatizaciones seleccionadas}
\label{vii:subsec:residual-aceleracion-ds}

El lector TRIT del estado \(x\) expresa \(s=\tau(x)\in\{+1,0,-1\}\),
y su reloj positivo expresa
\begin{equation}
 H_*=\frac{2\pi}{108t_0(x)}>0.
 \label{vii:eq:reloj-carta-ds}
\end{equation}
Sea \(h_s\) una métrica tridimensional de curvatura seccional constante
\(s\). Las tres coordenatizaciones de la realización seleccionada se escriben
\begin{equation}
 g_s=-\mathrm dt^2+a_s(t)^2h_s,\qquad
 a_+(t)=H_*^{-1}\cosh(H_*t),\quad
 a_0(t)=\exp(H_*t),\quad
 a_-(t)=H_*^{-1}\sinh(H_*t).
 \label{vii:eq:cartas-ds-residual}
\end{equation}
Las coordenatizaciones cerrada y plana se consideran para \(t\in\mathbb R\);
la coordenatización abierta, para \(t>0\). El signo \(s\) y el reloj \(H_*\)
mantienen sus antecedentes en el estado que selecciona la realización.

\begin{lemma}[Curvatura de las tres coordenatizaciones del mismo reloj]
\label{vii:lem:curvatura-cartas-ds}
Para las funciones de \eqref{vii:eq:cartas-ds-residual},
\begin{equation}
 \frac{\ddot a_s}{a_s}=H_*^2,\qquad
 \left(\frac{\dot a_s}{a_s}\right)^2+\frac{s}{a_s^2}=H_*^2,
 \qquad
 \operatorname{Ric}[g_s]=3H_*^2g_s.
 \label{vii:eq:curvatura-cartas-ds}
\end{equation}
En consecuencia, \(R[g_s]=12H_*^2\) y
\(\mathsf G[g_s]=-3H_*^2g_s\).
\end{lemma}
\begin{proof}
Las dos primeras igualdades se obtienen diferenciando las tres funciones:
en las coordenatizaciones cerrada y abierta se utiliza
\(\cosh^2z-\sinh^2z=1\), y en la plana
\(\dot a_0=H_*a_0\). Para verificar la identificación tensorial,
los coeficientes de conexión que involucran la coordenada temporal son
\[
 \Gamma^0{}_{ij}=a_s\dot a_s(h_s)_{ij},\qquad
 \Gamma^i{}_{0j}=(\dot a_s/a_s)\delta^i_j,
\]
y los puramente espaciales coinciden con los de \(h_s\).
La contracción de la curvatura da
\[
 \operatorname{Ric}_{00}=-3\ddot a_s/a_s,\qquad
 \operatorname{Ric}_{0i}=0,\qquad
 \operatorname{Ric}_{ij}
  =(a_s\ddot a_s+2\dot a_s^2+2s)(h_s)_{ij}.
\]
Sustituir las dos primeras identidades de
\eqref{vii:eq:curvatura-cartas-ds} produce
\(\operatorname{Ric}_{00}=-3H_*^2\) y
\(\operatorname{Ric}_{ij}=3H_*^2a_s^2(h_s)_{ij}\).
Tomar la traza y formar el tensor de Einstein prueba las otras
afirmaciones, con la convención de curvatura de
\eqref{vii:eq:ecuacion-metrica-efectiva}.
\end{proof}

\begin{theorem}[Densidad, presión y signo de la aceleración seleccionada]
\label{vii:thm:residual-ds-aceleracion}
En las coordenatizaciones anteriores, para \(\Lambda_{\rm ref}=0\) y \(T_b=0\),
el tensor residual es
\begin{equation}
 T^{\rm res}=-\rho_*g_s,\qquad
 \rho_*:=\frac{3H_*^2}{8\pi G_{\rm int}}>0,
 \qquad T^\Lambda=T^{\rm res},\quad T^0=0.
 \label{vii:eq:residual-ds-puro}
\end{equation}
Para un observador comóvil unitario \(u=\partial_t\), la densidad
y la presión isotrópica son \(\rho_*\) y \(p_*=-\rho_*\).
La ecuación de aceleración se evalúa como
\begin{equation}
 \frac{\ddot a_s}{a_s}
 =-\frac{4\pi G_{\rm int}}3(\rho_*+3p_*)=H_*^2>0.
 \label{vii:eq:aceleracion-residual-ds}
\end{equation}
\end{theorem}
\begin{proof}
El lema anterior y \eqref{vii:eq:residual-definicion} dan
\(T^{\rm res}=-(3H_*^2/\kappa_x)g_s\).
Como \(g_s(u,u)=-1\),
\(T^{\rm res}_{\mu\nu}u^\mu u^\nu=\rho_*\).
El proyector espacial \(q^{\mu\nu}=g_s^{\mu\nu}+u^\mu u^\nu\)
satisface \(q^{\mu\nu}(g_s)_{\mu\nu}=3\), y por ello
\[
 p_*:=\tfrac13q^{\mu\nu}T^{\rm res}_{\mu\nu}=-\rho_*.
\]
La traza vale \(-4\rho_*\), de modo que
\eqref{vii:eq:descomposicion-traza-residual} produce los dos sectores
de \eqref{vii:eq:residual-ds-puro}. Finalmente,
\(\rho_*+3p_*=-2\rho_*\), y la sustitución en el segundo miembro
de \eqref{vii:eq:aceleracion-residual-ds} da
\(\kappa_x\rho_*/3=H_*^2\), igual al valor geométrico ya calculado.
La positividad utiliza exactamente \(G_{\rm int}>0\) y \(H_*>0\).
\end{proof}

En esta realización, el sector de traza expresa una densidad constante
y una presión opuesta, y el reloj interno fija la aceleración positiva.
En una coordenatización general, \(T^\Lambda\) describe la componente isotrópica
puntual, mientras \(T^0\) conserva el residual de traza nula y ambos
obedecen el intercambio de \eqref{vii:eq:balances-sectores-residuales}.
La identificación de estos sectores con fuentes oscuras observacionales
requiere que una misma realización reproduzca conjuntamente expansión,
lentes, rotación, crecimiento de estructura y fondo cósmico. La
construcción presente fija el tensor, sus balances y la aceleración de
las coordenatizaciones seleccionadas que intervienen en esa comparación.
